\subsection{Widlar Mirror}

\subsubsection{Overview}
Constructed from two transistors, It is the
simplest form of a current mirror that can provide a
current gain different from unity.

\subsubsection{Schematic}
\begin{circuitfig}
  \draw (2.98, 7.905) to[american voltage source, l_={$V_{cc}$}] (1.73, 7.905);
	\node[nmos, arrowmos, xscale=-1] (N1) at (3.25, 4.405){} node[anchor=east] at (N1.text){$M_0$};
	\node[nmos, arrowmos] (N2) at (5.75, 4.405){} node[anchor=west] at (N2.text){$M_1$};
	\draw (3.25, 7.405) to[european resistor, l={$G_2$}] (3.25, 5.905);
	\draw (5.75, 5.655) -| (5.75, 5.175);
	\draw (4.5, 4.405) |- (3.25, 5.405);
	\draw (4.77, 4.405) -- (4.23, 4.405);
	\draw (3.25, 3.635) -| (3.25, 3.405) -| (5.75, 3.635);
	\draw (3.75, 2.655) to[american voltage source, l={$V_{ee}$}] (2.5, 2.655);
	\draw (3.75, 2.655) -| (4.5, 3.405);
	\node[ground, rotate=-90] at (1.73, 7.905){};
	\node[ground, rotate=-90] at (2.5, 2.655){};
	\draw[dash pattern={on 1.6pt off 0.8pt}, -stealth] (5.75, 5.655) -- (6.25, 5.655);
	\node[shape=rectangle, minimum width=0.965cm, minimum height=0.465cm] (N3) at (6.75, 5.615){} node[anchor=center] at (N3.text){$V_{out}$};
	\draw (5.75, 7.365) to[european resistor, l={\textcolor{rgb,255:red,255;green,0;blue,0}{$G_3$}}] (5.75, 5.865);
	\draw (5.75, 5.655) -- (5.75, 5.885);
	\draw (3.25, 7.405) |- (2.98, 7.905);
	\draw (3.25, 7.905) -| (5.75, 7.365);
	\node[shape=rectangle, minimum width=4.965cm, minimum height=1.09cm] at (4.375, 9.928){} node[anchor=north, align=center, text width=4.577cm, inner sep=6pt] at (4.375, 10.49){Large Signal Schematic\\(Sink)};
	\node[shape=rectangle, draw={rgb,255:red,255;green,0;blue,0}, line width=1pt, dash pattern={on 4pt off 2pt on 1pt off 2pt on 1pt off 2pt}, minimum width=1.465cm, minimum height=2.195cm] at (6, 7){};
	\node[shape=rectangle, minimum width=3.09cm, minimum height=2.465cm] at (8.438, 6.865){} node[anchor=north west, align=left, text width=2.702cm, inner sep=6pt] at (6.875, 8.115){\textcolor{rgb,255:red,255;green,0;blue,0}{\footnotesize Exists Only For Analysis. In  Implementation take output directly from V\_\{out\}}};
	\node[shape=rectangle, minimum width=0.965cm, minimum height=0.465cm] (N4) at (2.25, 5.405){} node[anchor=center] at (N4.text){$V_{x}$};
	\draw[dash pattern={on 1.6pt off 0.8pt}, -stealth] (3.25, 5.405) |- (2.75, 5.405);
	\node[shape=rectangle, minimum width=0.965cm, minimum height=0.465cm] (N5) at (4.5, 6.115){} node[anchor=center] at (N5.text){$I_{G}$};
	\draw (3.25, 5.175) -| (3.25, 5.405);
	\draw (3.25, 5.405) -| (3.25, 5.905);
	\draw (10, 1.75) |- (10, 11);
	\draw (17.25, 5.563) to[american current source, l_={$g_{m_1}v_{gs}$}] (17.25, 4.063);
	\draw (18.75, 5.563) to[european resistor, l_={$g_{o_1}$}] (18.75, 4.063);
	\draw (17.25, 4.063) -- (18.75, 4.063);
	\draw (18.75, 5.563) -- (17.25, 5.563);
	\draw (17.25, 4.063) -- (16, 4.063);
	\draw (16, 5.563) -- (15.25, 5.563);
	\node[shape=rectangle, minimum width=0.465cm, minimum height=0.465cm] at (16, 5.313){} node[anchor=center, align=center, text width=0.077cm, inner sep=6pt] at (16, 5.313){\footnotesize G};
	\node[shape=rectangle, minimum width=0.465cm, minimum height=0.465cm] at (16, 4.313){} node[anchor=center, align=center, text width=0.077cm, inner sep=6pt] at (16, 4.313){\footnotesize S};
	\node[shape=rectangle, minimum width=0.465cm, minimum height=0.465cm] at (19.25, 5.313){} node[anchor=center, align=center, text width=0.077cm, inner sep=6pt] at (19.25, 5.313){\footnotesize D};
	\draw (15.25, 5.563) -- (14.5, 5.563);
	\draw (13, 5.563) to[american current source, l={$g_{m_0}v_{gs}$}] (13, 4.063);
	\draw (11.5, 5.563) to[european resistor, l={$g_{o_0}$}] (11.5, 4.063);
	\draw (13, 5.563) -- (11.5, 5.563);
	\draw (13, 4.063) -- (11.5, 4.063);
	\draw (14.5, 4.063) -- (13, 4.063);
	\node[shape=rectangle, minimum width=0.465cm, minimum height=0.465cm] at (14.5, 5.313){} node[anchor=center, align=center, text width=0.077cm, inner sep=6pt] at (14.5, 5.313){\footnotesize G};
	\node[shape=rectangle, minimum width=0.465cm, minimum height=0.465cm] at (14.5, 4.313){} node[anchor=center, align=center, text width=0.077cm, inner sep=6pt] at (14.5, 4.313){\footnotesize S};
	\node[shape=rectangle, minimum width=0.465cm, minimum height=0.465cm] at (11, 5.313){} node[anchor=center, align=center, text width=0.077cm, inner sep=6pt] at (11, 5.313){\footnotesize D};
	\draw (11.5, 7.563) to[european resistor, l={$G_2$}] (11.5, 6.063);
	\draw (11.5, 6.063) -| (11.5, 5.562);
	\draw (11.5, 7.563) -| (11.5, 8.062) -- (12.25, 8.062) -| (12.25, 7.812);
	\draw (18.75, 7.563) to[european resistor, l={\textcolor{rgb,255:red,255;green,0;blue,0}{$G_3$}}] (18.75, 6.063);
	\draw (18.75, 6.063) -| (18.75, 5.562);
	\draw (18.75, 7.563) -| (18.75, 8.062) |- (18, 8.062) -| (18, 7.812);
	\draw[dash pattern={on 1.6pt off 0.8pt}, -stealth] (18.75, 5.563) -- (19.25, 5.563);
	\node[shape=rectangle, minimum width=0.715cm, minimum height=0.465cm] (N6) at (10.625, 5.563){} node[anchor=center] at (N6.text){$V_{x}$};
	\node[shape=rectangle, minimum width=0.965cm, minimum height=0.465cm] (N7) at (19.75, 5.563){} node[anchor=center] at (N7.text){$V_{out}$};
	\node[ground] at (14.5, 4.063){};
	\node[ground] at (16, 4.063){};
	\node[ground] at (12.25, 7.812){};
	\node[ground] at (18, 7.812){};
	\node[shape=rectangle, minimum width=0.965cm, minimum height=0.465cm] (N8) at (15.25, 6.813){} node[anchor=center] at (N8.text){$I_{G}$};
	\draw[dash pattern={on 1.6pt off 0.8pt}, -stealth] (11.5, 5.563) -- (11, 5.563);
	\node[shape=rectangle, minimum width=4.215cm, minimum height=1.09cm] at (15.375, 9.937){} node[anchor=north, align=center, text width=3.827cm, inner sep=6pt] at (15.375, 10.5){Small Signal Schematic\\(Sink)};
	\node[circ] at (3.25, 5.405){};
	\node[circ] at (4.5, 4.405){};
	\node[circ] at (4.5, 3.405){};
	\node[circ] at (3.25, 7.905){};
	\draw (13, 5.563) |- (15.25, 6.063) -| (15.25, 5.563);
	\draw[dash pattern={on 1.6pt off 0.8pt}, -stealth] (15.25, 6.063) -- (15.25, 6.563);
	\draw[dash pattern={on 1.6pt off 0.8pt}, -stealth] (4.5, 5.365) -- (4.5, 5.865);
	\node[circ] at (11.5, 5.563){};
	\node[circ] at (15.25, 5.563){};
	\node[circ] at (18.75, 5.563){};
	\node[circ] at (13, 5.563){};
	\node[shape=rectangle, draw={rgb,255:red,255;green,0;blue,0}, line width=1pt, dash pattern={on 4pt off 2pt on 1pt off 2pt on 1pt off 2pt}, minimum width=1.465cm, minimum height=1.965cm] at (19, 7.063){};
\end{circuitfig}

\subsubsection{Transfer Function}
Here's a table of the voltages at each terminal of the transistors:

\begin{tabularx}{\linewidth}{|X|X|X|X|X|X|} \hline
  Transistor & Gate (G) & Source (S) & Drain (D) & $V_{gs}$ & $V_{ds}$ \\ \hline
  $M_0$ & $V_{x}$ & $V_{ee}$ & $V_{x}$ & $V_{x} - V_{ee}$ & $V_{x} - V_{ee}$ \\ \hline
  $M_1$ & $V_{x}$ & $V_{ee}$ & $V_{out}$ & $V_{x} - V_{ee}$ & $V_{out} - V_{ee}$ \\ \hline
\end{tabularx}

\begin{enumerate}
  \item \textbf{Large Signal Analysis}:\\
    Applying KCL at $V_{out}$:
    \begin{align*}
      (V_{out} - V_{cc})G_3 + I_{M_1} &= 0 \\
      V_{out}G_3 - V_{cc}G_3 + I_{M_1} &= 0 \\
      V_{out} &= V_{cc} - \frac{I_{M_1}}{G_3} \tag{1}\end{align*}
    Applying KCL at $V_{x}$:
    \begin{align*}
      (V_{x} - V_{cc})G_2 + I_{M_0} + I_{G} &= 0 \\
      V_{x}G_2 - V_{cc}G_2 + (I_{M_0} + I_{G}) &= 0 \\
      V_{x}G_2 &= - (I_{M_0} + I_{G}) + V_{cc}G_2 \\
      V_{x} &= V_{cc} - \frac{(I_{M_0} + I_{G})}{G_2} \tag{2}
    \end{align*}
    Using the transistor equations for $I_{M_0}$ and
    $I_{M_1}$:
    \begin{align*}
      I_{M_0} &= k_{n_0}(V_{gs_0} - V_{th})^2 \cdot (1 + \lambda V_{ds_0}) \\
              &= k_{n_0}(V_{x} - V_{ee} - V_{th})^2 \cdot (1 + \lambda (V_{x} - V_{ee})) \\
              &= k_{n_0}\left(\left(V_{cc} - \frac{(I_{M_0} + I_{G})}{G_2}\right) - V_{ee} - V_{th}\right)^2 \cdot \left(1 + \lambda \left(\left(V_{cc} - \frac{(I_{M_0} + I_{G})}{G_2}\right) - V_{ee}\right)\right) \\
              &= k_{n_0}\left(\left(V_{cc} - \frac{(I_{M_0} + I_{G})}{G_2}\right) - V_{ee} - V_{th}\right)^2 \quad \text{[assuming } \lambda = 0 \text{]} \tag{3} \\ \\
      I_{M_1} &= k_{n_1}(V_{gs_1} - V_{th})^2 \cdot (1 + \lambda V_{ds_1}) \\
              &= k_{n_1}(V_{x} - V_{ee} - V_{th})^2 \cdot (1 + \lambda (V_{out} - V_{ee})) \\
              &= k_{n_1}\left(\left(V_{cc} - \frac{(I_{M_0} + I_{G})}{G_2}\right) - V_{ee} - V_{th}\right)^2 \cdot \left(1 + \lambda \left(V_{cc} - \frac{I_{M_1}}{G_3} - V_{ee}\right)\right) \\
              &= k_{n_1}\left(\left(V_{cc} - \frac{(I_{M_0} + I_{G})}{G_2}\right) - V_{ee} - V_{th}\right)^2 \quad \text{[assuming } \lambda = 0 \text{]} \tag{4} \\
    \end{align*}
    It can be seen that,
    \begin{align*}
      I_{in}  = I_{M_0} &\approx I_{M_1}  \quad \text{[if } k_{n_0} = k_{n_1} \text{]} \\
      I_{out} = I_{M_1}
    \end{align*}
    Thus,
    \begin{align*}
      \frac{I_{out}}{I_{in}} \approx \frac{I_{M_1}}{I_{M_0}} \approx \frac{k_{n_1}}{k_{n_0}} = \frac{\frac{\mu_n C_{ox} W_1}{L_1}}{\frac{\mu_n C_{ox} W_0}{L_0}} = \frac{\frac{W_1}{L_1}}{\frac{W_0}{L_0}} = \frac{W_1 L_0}{L_1 W_0}
    \end{align*}
  \item \textbf{Small Signal Analysis}:\\
    Applying KCL at $V_{out}$:
    \begin{align*}
      V_{out}G_3 + V_{out}g_{o_1} + g_{m_1}V_{x} &= 0 \\
      V_{out}(G_3 + g_{o_1}) + g_{m_1}V_{x} &= 0 \tag{5}
    \end{align*}
    Applying KCL at $V_{x}$:
    \begin{align*}
      V_{x}G_2 + V_{x}g_{o_0} + g_{m_0}V_{x} + I_{G} &= 0 \\
      V_{x}(G_2 + g_{o_0} + g_{m_0}) &= 0 \quad \text{[assuming } I_G \approx 0 \text{ by MOS physics]} \tag{6}
    \end{align*}
    Adding voltage source $V_s$ at $V_x$ that supplies
    a current $I_s$ into $V_{x}$, we get:
    \begin{align*}
      V_{x}(G_2 + g_{o_0} + g_{m_0}) - I_{s} &=0 \tag{7} \\
      V_{x} = V_{s} \tag{8}
    \end{align*}
    Solving equations (5), (7) and (8):
    \begin{align*}
      V_{out} &= - \frac{g_{m_1} V_{x}}{G_3 + g_{o_1}} \\
              &= - \frac{g_{m_1} V_{s}}{G_3 + g_{o_1}} \quad \text{[from (8)]} \\
      H(s) = \frac{N(s)}{D(s)} = \frac{V_{out}}{V_{s}} = - \frac{g_{m_1}}{G_3 + g_{o_1}}
    \end{align*}
\end{enumerate}

\subsubsection{Analysis}
\begin{enumerate}
  \item \textbf{Output Load (Max)}: Maximum possible
    load that can be connected at the output before
    the transistor $M_2$ goes out of saturation can
    be calculated as follows:
    \begin{align*}
      R_{out, \max} &= \frac{V_{L, \max}}{I_{M_1}} \quad \text{[where, } V_{ee} \leq V_{L, \max} \leq V_{cc} \text{]}
    \end{align*}
  \item \textbf{Impedances}:
    \paragraph{At Input:} Applying a test current $I_{test}$
      at $V_{x}$, from equation (6):
      \begin{align*}
        V_{x}(G_2 + g_{o_0} + g_{m_0}) &= I_{test} \\
        \therefore Z_{in} &= \frac{V_{x}}{I_{test}} = \frac{1}{G_2 + g_{o_0} + g_{m_0}}
      \end{align*}
      Ignoring $G_2$ and assuming $g_{m_0} \gg g_{o_0}$, we get:
      \begin{align*}
        Z_{in} &\approx \frac{1}{g_{m_0}} = \frac{2 I_{M_0}}{V_{ov_0}}
      \end{align*}
    \paragraph{At Output:} Applying a test current $I_{test}$
      at $V_{out}$, from equation (5):
      \begin{align*}
        V_{out}(G_3 + g_{o_1}) + g_{m_1}V_{x} &= I_{test} \\
        V_{out}(G_3 + g_{o_1}) + g_{m_1}(0) &= I_{test} \quad \text{[from (7)]} \\
        \therefore Z_{out} &= \frac{V_{out}}{I_{test}} = \frac{1}{G_3 + g_{o_1}}
      \end{align*}
      Ignoring $G_3$, we get:
      \begin{align*}
        Z_{out} &\approx \frac{1}{g_{o_1}} = \frac{1}{\lambda_{M_1} I_{M_1}}
      \end{align*}
\end{enumerate}

\subsubsection{Design}
Large Signal Analysis can be used to set the bias
currents and voltages in this circuit.
\begin{enumerate}
  \item Choose appropriate transistor dimensions
    ($W$ and $L$) for transistors $M_0$ and $M_1$
    to achieve the desired current gain.
    \begin{align*}
      \frac{I_{out}}{I_{in}} &= \frac{W_1 L_0}{L_1 W_0}
    \end{align*}
  \item Choose a suitable overdrive voltage value
    ($V_{ov}$) for computing the bias voltage at
    node $V_{x}$.
    \begin{align*}
      V_{ov_0} &= V_{gs_0} - V_{th} \\
      V_{ov_0} &= (V_{x} - V_{ee}) - V_{th} \\
      V_{x} &= V_{ov_0} + V_{th} + V_{ee}
    \end{align*}
  \item Using the chosen $V_{ov}$, compute $W_0$
    to set the required bias current $I_{in}$.
    \begin{align*}
      I_{in} = I_{M_0} &= \frac{1}{2} \mu_n C_{ox} \frac{W_0}{L_0} V_{ov_0}^2 \\
      \therefore W_0 &= \frac{2 I_{in} L_0}{\mu_n C_{ox} V_{ov_0}^2}
    \end{align*}
    $W_1$ can then be computed using the current
    gain relation from step (a).
  \item Compute the required value of $G_2$ (or
    $R_2$) to ensure that the voltage at node
    $V_{x}$ is as calculated in step (b).
    \begin{align*}
      (V_{x} - V_{cc})G_2 + I_{M_0} + I_{G} &= 0 \\
      (V_{x} - V_{cc})G_2 + I_{M_0} &\approx 0 \quad \text{[assuming } I_G \approx 0 \text{ by MOS physics ]} \\
      (V_{x} - V_{cc})G_2 &\approx - I_{M_0} \\
      G_2 &\approx - \frac{I_{M_0}}{(V_{x} - V_{cc})} \\
      \therefore G_2 &= \frac{I_{M_0}}{V_{cc} - V_{x}} \\
      R_2 &= \frac{1}{G_2} = \frac{V_{cc} - V_{x}}{I_{M_0}}
    \end{align*}
\end{enumerate}`

\subsubsection{Question}
Design a Widlar Current Mirror (Sink) that meets
the following requirements. The mirror must provide
a current gain at the output of 5. The input current
must be $I_{in} = 100 \mu A$. The mirror must operate
with $V_{cc} = 5V$ and $V_{ee} = 0V$. Assume the
following transistor parameters: $V_{th} = 1V$,
$\lambda = 0.02 V^{-1}$, $\frac{1}{2} \mu_n C_{ox} =
100 \mu A/V^2$ and $L = 1 \mu m$ for all transistors.
Find the output impedance of the mirror and the
maximum load that can be connected at the output
if $V_{L, \max} = 4V$.
\paragraph{Solution:}
\begin{enumerate}
  \item \textbf{Choose appropriate transistor dimensions:}
    \begin{align*}
      \frac{I_{out}}{I_{in}} &= \frac{W_1 L_0}{L_1 W_0} = 5 \\
      \therefore W_1 &= 5 W_0 \quad \text{[assuming } L_0 = L_1 \text{]}
    \end{align*}
  \item \textbf{Choose a suitable overdrive voltage and determine
    the resulting bias voltages:}
    \begin{align*}
      V_{x} &= V_{ov_0} + V_{th} + V_{ee} = 0.2V + 1V + 0V = 1.2V
    \end{align*}
  \item \textbf{Compute $W_0$ to set the required bias current:} For $I_{in} = 100 \mu A$
    \begin{align*}
      W_0 &= \frac{2 I_{in} L_0}{\mu_n C_{ox} V_{ov_0}^2} \\
          &= \frac{2 \cdot 100 \mu A \cdot 1 \mu m}{100 \mu A/V^2 \cdot (0.2V)^2} \\
          &= 25 \mu m \\
      \therefore W_1 &= 5 W_0 = 125 \mu m
    \end{align*}
  \item \textbf{Compute the required value of $G_2$ (or $R_2$):}
    \begin{align*}
      R_2 &= \frac{1}{G_2} = \frac{V_{cc} - V_{x}}{I_{M_0}} \\
          &= \frac{5V - 1.2V}{100 \mu A} \\
          &= 38 k\Omega
    \end{align*}
  \item \textbf{Compute the output impedance:}
    \begin{align*}
      Z_{out} &\approx \frac{1}{g_{o_1}} \approx \frac{1}{\lambda_{M_1} I_{M_1}} \\
              &= \frac{1}{0.02 V^{-1} \cdot 500 \mu A} \\
              &= 100 k\Omega
    \end{align*}
  \item \textbf{Compute the maximum load that can be connected
    at the output:}
    \begin{align*}
      R_{out, \max} &= \frac{V_{L, \max}}{I_{M_1}} \\
                   &= \frac{4V}{500 \mu A} \\
                   &= 8 k\Omega
    \end{align*}
\end{enumerate}

\subsubsection{Code}
\begin{lstlisting}[language=Octave]
  clear; clc;
  pkg load symbolic;

  syms Vx Vout Vs Is Itest
  syms G3 G2 go0 gm0 go1 gm1

  % system kcl
  eq1 = Vout*(G3 + go1) + Vx*gm1 == 0; % KCL @ Vout
  eq2 = Vx*(G2 + go0 + gm0) == 0; % KCL @ Vx

  % voltage transfer function: Vs forced at Vx & Is injected into Vx
  eq2_vt = Vx*(G2 + go0 + gm0) - Is == 0;
  eq3_vt = Vx == Vs; % equation introduced.
  sol_vt = solve([eq1, eq2_vt, eq3_vt], [Vout, Vx, Vs]);
  volt_transfer = simplify(sol_vt.Vout/sol_vt.Vs);
  % volt_transfer = subs(volt_transfer, 'G2', 0);
  % volt_transfer = subs(volt_transfer, 'G3', 0);
  printf('=== Voltage Transfer Results ===\n');
  pretty(volt_transfer)
  printf('\n');

  % current transfer function: Is injected into Vx
  sol_ct = solve([eq1, eq2_vt], [Vout, Vx]);
  Iout = sol_ct.Vout * G3;
  curr_transfer = simplify(Iout/Is);
  printf('=== Current Transfer Results ===\n');
  pretty(curr_transfer)
  printf('\n');

  % input impedance: Itest injected into Vx
  eq2_zin = Vx*(G2 + go0 + gm0) + Itest == 0;
  sol_zin = solve([eq1, eq2_zin], [Vout, Vx]);
  zin = simplify(sol_zin.Vx/Itest);
  % zin = subs(zin, 'G2', 0);
  % zin = subs(zin, 'G3', 0);
  printf('=== Input Impedance Results ===\n');
  pretty(zin)
  printf('\n');

  % output impedance: Itest injected into Vout
  eq1_zout = Vout*(G3 + go1) + Vx*gm1 + Itest == 0;
  sol_zout = solve([eq1_zout, eq2], [Vout, Vx]);
  zout = simplify(sol_zout.Vout/Itest);
  % zout = subs(zout, 'G2', 0);
  % zout = subs(zout, 'G3', 0);
  printf('=== Output Impedance Results ===\n');
  pretty(zout)
  printf('\n');
\end{lstlisting}
